\subsubsection{Annotation Operations}
\label{sec:geo_annotation_ops}  

Annotations are listed in the 'Annotations' node of the 'Layers' tab. \\

There are two kinds of annotations listed under this node.

\paragraph*{View Point} Annotations provided by an active View Point, e.g. selected evaluation results or externally provided annotation data.
See \nameref{sec:key_concepts_viewpoints} for an introductory explanation about View Points. \\

View Point annotations persist as long as the providing View Point is active. \\

Clicking on the symbol of the 'View Point' node will open a context menu containing the following operations:

\begin{table}[H]
    \center
    \begin{tabular}{ | l | l |}
      \hline
      \textbf{Operation} & \textbf{Description} \\ \hline
      Show & Shows all annotations of a selected name \\ \hline
      Hide & Hides all annotations of a selected name \\ \hline
      Remove All & Removes the View Point annotations from the view \\ \hline
    \end{tabular}
\end{table}

The 'Show' and 'Hide' submenus list only annotation names that exist more than once. 'Remove All' is only listed if View Point annotations exist, and removes them in this Geographic View only. \\

\begin{figure}[H]
  \center
    \fbox{\parbox{0.8\textwidth}{\vspace{1.5cm}\centering \textcolor{red}{\textbf{[FIGURE TODO: RMB click on the symbol of the 'View Point' node in the 'Layers' tab, with a View Point active that provides several annotations of the same name]}}\vspace{1.5cm}}}
  \caption{Geographic View 'View Point' node operations}
\end{figure}

\paragraph*{Grids} Grid data imported into Geographic View. Grid data can be imported into 
Geographic View either by sending it from a Grid View (see \nameref{sec:grid_export_geographic}), or as an external GeoTIFF file (see \nameref{sec:geo_geotiff_import}). \\

Imported grid items persist in a Geographic View until COMPASS is closed. \\

Clicking on the symbol of the 'Grids' node will open a context menu containing the following operations:

\begin{table}[H]
    \center
    \begin{tabular}{ | l | l |}
      \hline
      \textbf{Operation} & \textbf{Description} \\ \hline
      Remove All & Removes all imported grids from the view \\ \hline
    \end{tabular}
\end{table}

\begin{figure}[H]
    \centering
    \hspace*{-2.5cm}
    \includegraphics[width=6cm,frame]{figures/geoview_gridops_grids.png}
  \caption{Geographic View 'Grids' node operations}
\end{figure}

Clicking on the symbol of an imported grid item will open a context menu containing the following operations:

\begin{table}[H]
    \center
    \begin{tabular}{ | l | l |}
      \hline
      \textbf{Operation} & \textbf{Description} \\ \hline
      Zoom to Item & Zooms to the grid item \\ \hline
      Rename & Renames the grid item to a specified name \\ \hline
      Export & Exports the grid item to a GeoTIFF file \\ \hline
      Opacity & Sets the opacity of a grid item \\ \hline
      Remove & Removes the grid item from the view \\ \hline
    \end{tabular}
\end{table}

\begin{figure}[H]
    \centering
    \hspace*{-2.5cm}
    \includegraphics[width=6cm,frame]{figures/geoview_gridops_grid.png}
  \caption{Geographic View grid item operations}
\end{figure}
\ \\

All annotation items can be shown or hidden by using the respective check box next to them.
